\documentclass[main.tex]{subfiles}
\begin{document}

\section{Exercise 2 – NMOS Common-Source Stage with Diode-Connected Load}

\subsection*{Given}
\[
\begin{aligned}
\mu_n C_{ox} &= 134~\mu\text{A/V}^2, &
(W/L)_1 &= 100, &
(W/L)_2 &= 20,\\
V_T &= 0.7~\text{V}, &
\lambda &= 0.1~\text{V}^{-1}, &
I_D &= 0.58~\text{mA}, &
V_{DD} &= 3.04~\text{V}.
\end{aligned}
\]

\subsection*{Method}
For each NMOS transistor in saturation:
\begin{equation}
I_D = \tfrac{1}{2}\mu_n C_{ox}\tfrac{W}{L}(V_{GS}-V_T)^2(1+\lambda V_{DS}),
\label{eq:2.1}
\end{equation}
and the effective voltage:
\begin{equation}
V_{\text{eff}} = V_{GS}-V_T = \sqrt{\tfrac{2I_D}{\mu_n C_{ox}(W/L)}}.
\label{eq:2.2}
\end{equation}

For small-signal operation:
\begin{equation}
g_m=\tfrac{2I_D}{V_{\text{eff}}}, \qquad
r_o=\tfrac{1}{\lambda I_D}.
\label{eq:2.3}
\end{equation}

Because \(M_2\) is diode-connected, its small-signal resistance is approximately
\[
r_{d2}=\frac{1}{g_{m2}+g_{ds2}}\approx\frac{1}{g_{m2}},
\]
so the load seen by \(M_1\) is dominated by \(1/g_{m2}\).
The small-signal voltage gain becomes
\begin{equation}
A_v \approx -g_{m1}\Big(\frac{1}{g_{m2}}\Big) = -\frac{g_{m1}}{g_{m2}}.
\label{eq:2.4}
\end{equation}

\subsection*{Calculation}
\[
\begin{aligned}
V_{\text{eff1}} &= 0.294~\text{V}, & V_{GS1}&=0.994~\text{V},\\
V_{\text{eff2}} &= 0.657~\text{V}, & V_{GS2}&=1.357~\text{V}.
\end{aligned}
\]
Since \(M_2\) is diode-connected:
\[
V_{out}=V_{DD}-V_{GS2}=1.68~\text{V}.
\]
\[
g_{m1}=\frac{2I_D}{V_{\text{eff1}}}=\frac{2(0.58~\text{mA})}{0.294}=3.95~\text{mS},\qquad
g_{m2}=\frac{2I_D}{V_{\text{eff2}}}=\frac{2(0.58~\text{mA})}{0.657}=1.77~\text{mS}.
\]
\[
\boxed{A_v = -\frac{g_{m1}}{g_{m2}} = -\frac{3.95}{1.77} = -2.23~\text{V/V}.}
\]

\subsection*{Result}
\[
\boxed{V_{out}=1.68~\text{V},\quad A_v\approx-2.2~(6.9~\text{dB}).}
\]
\subsection*{Region check}
\[
V_{DS1}=1.68~\text{V}>V_{\text{eff1}}=0.29~\text{V},\quad
V_{DS2}=V_{GS2}=1.36~\text{V}>V_{\text{eff2}}=0.66~\text{V}.
\]
\[
\boxed{\text{Both transistors are in saturation.}}
\]

\end{document}
